面向多计算管道和共享外存的神经网络处理器，本文在算子依赖约束下联合选择子图划分、核心映射与执行顺序，以官方仿真总完成时间为主要目标，同时间下优先减少额外数据搬运。模型按跨子图经DDR、同核驻留与跨核同步、共享只读FIFO缓存三种机制逐层扩展。

对问题一，采用无环有向图划分与多核列表调度，从完整分量、共享输入与有界局部切分生成候选，保护基础划分机会并进行关键区域联合调整。百图五核逐图平均加速比为3.665903，总周期为73170501。

对问题二，按跨核搬运完成后500周期释放建立事件约束，采用保护基础搜索的结构算法组合，并复用不变的候选生成准备。五核逐图平均加速比为4.264407，总周期为65342192。对问题三，加入共享FIFO状态和DDR、L2独立带宽池，在普通结构候选之外引入缓存事件邻域；五核平均加速比为4.364503，总周期为63759268。固定最终方案、仅切换L2时，相同核数的平均时间比为1.030244，与跨方案综合比1.027332分开报告。

三问均保存100图、1至5核的逐配置官方结果。结果表明，减少搬运或增加缓存命中并不保证缩短关键路径，应以完整事件仿真裁决。算法不穷举组合空间；冷启动大图压力核验与历史初解成本单列，尚不能把限时搜索结果解释为所有图端到端均在十分钟内完成，也未证明全局最优。
